%%%

\documentclass[letterpaper, 10 pt, conference]{ieeeconf}  % Comment this line out if you need a4paper

%\documentclass[a4paper, 10pt, conference]{ieeeconf}      % Use this line for a4 paper

\IEEEoverridecommandlockouts                              % This command is only needed if 
                                                          % you want to use the \thanks command

\overrideIEEEmargins                                      % Needed to meet printer requirements.

%In case you encounter the following error:
%Error 1010 The PDF file may be corrupt (unable to open PDF file) OR
%Error 1000 An error occurred while parsing a contents stream. Unable to analyze the PDF file.
%This is a known problem with pdfLaTeX conversion filter. The file cannot be opened with acrobat reader
%Please use one of the alternatives below to circumvent this error by uncommenting one or the other
%\pdfobjcompresslevel=0
%\pdfminorversion=4

% See the \addtolength command later in the file to balance the column lengths
% on the last page of the document

% The following packages can be found on http:\\www.ctan.org
%\usepackage{graphics} % for pdf, bitmapped graphics files
%\usepackage{epsfig} % for postscript graphics files
%\usepackage{mathptmx} % assumes new font selection scheme installed
%\usepackage{times} % assumes new font selection scheme installed
%\usepackage{amsmath} % assumes amsmath package installed
%\usepackage{amssymb}  % assumes amsmath package installed

\usepackage{amsmath}
\usepackage{amssymb}

\usepackage{graphicx}

\usepackage{subcaption}
\usepackage{balance}
\usepackage{dblfloatfix}

% table
\usepackage{booktabs}
\usepackage{multirow}

% \title{MeCaRL: Memory-enhanced Collision-aware Reinforcement Learning for UAV Navigation in Multi-scale Obstacle Environments}
% \title{\LARGE \bf
% Preparation of Papers for IEEE Sponsored Conferences \& Symposia*
% }

\title{\LARGE \bf
MeCaRL: Memory-enhanced Collision-aware Reinforcement Learning for UAV Navigation in Multi-scale Obstacle Environments
}

\author{Hong Hong, Feiyu Liao, Haitao Wang, and Hejun Wu%
\thanks{Hong Hong and Feiyu Liao contributed equally to this work. (Corresponding author: Haitao Wang.)}%
\thanks{All authors are with the School of Computer Science and Engineering, Sun Yat-sen University, Guangzhou, China, 
and the Guangdong Key Laboratory of Big Data Analysis and Processing, Guangzhou, China (e-mail: 
hongh2688@mail2.sysu.edu.cn; feiyuliao@mail2.sysu.edu.cn; wanght39@mail2.sysu.edu.cn; wuhejun@mail.sysu.edu.cn).}%
\thanks{Website with code: https://hongh26.github.io/mecarl/}%
}


\begin{document}



\maketitle
\thispagestyle{empty}
\pagestyle{empty}


%%%%%%%%%%%%%%%%%%%%%%%%%%%%%%%%%%%%%%%%%%%%%%%%%%%%%%%%%%%%%%%%%%%%%%%%%%%%%%%%
% \begin{abstract}
% Autonomous flight in unknown and cluttered environments remains a current research hotspot. Existing methods, when constructing experimental scenarios, mainly focus on the number of obstacles and their spatial density, while paying less attention to the perception and decision-making challenges caused by variations in obstacle size. To the best of our knowledge, we are among the first to introduce multi-scale obstacle scenarios into learning-based UAV visual navigation and proposes a memory-enhanced collision-aware reinforcement learning framework, MeCaRL. MeCaRL learns a collision-risk-oriented depth representation to improve sensitivity to small-sized obstacles, and incorporates a temporal memory mechanism to fuse historical observations and alleviate occlusions caused by large-sized obstacles. These two components complement each other and jointly improve the effectiveness of obstacle avoidance decisions. Experimental results demonstrate that MeCaRL not only performs well in nominal-size obstacle scenes, but also achieves safe and efficient autonomous flight in cluttered environments containing both large-sized and small-sized obstacles. Finally, real-world UAV experiments further verify the effectiveness of the proposed method in physical environments.

% \end{abstract}

\begin{abstract}
Unmanned aerial vehicle (UAV) navigation in unknown and cluttered environments remains an important research topic. 
Existing learning-based methods typically evaluate performance by varying obstacle number and density, while paying less attention to the perception and decision-making difficulties introduced by obstacle scale variations. 
In this work, we propose MeCaRL, a memory-enhanced and collision-aware reinforcement learning framework for UAV visual navigation. 
MeCaRL learns a collision-aware latent representation from depth observations to improve sensitivity to small obstacles, and incorporates temporal memory to fuse historical observations and mitigate occlusions caused by large obstacles. 
We evaluate MeCaRL in multi-scale obstacle environments, including ultra-small and extra-large obstacle settings. 
Results show that MeCaRL consistently outperforms strong learning-based baselines in success rate. 
The gains are particularly clear in ultra-small and extra-large scenarios, while performance remains strong in nominal-scale environments. 
Real-world outdoor flight experiments further demonstrate the feasibility of deploying MeCaRL for safe and efficient navigation in cluttered environments.
% Website with code: https://hongh26.github.io/mecarl/
\end{abstract}


\begin{figure}[t]
  \centering
  \includegraphics[scale=0.25]{figures/fig1_motivation.pdf}
  \caption{(a) Representative trajectories of Agile-autonomy, MAVRL, and MeCaRL (ours) in extra-large (left) and ultra-small (right) obstacle settings. (b) Success rates of the three methods in the two scenarios.}
  \label{fig:1}
\end{figure}

%%%%%%%%%%%%%%%%%%%%%%%%%%%%%%%%%%%%%%%%%%%%%%%%%%%%%%%%%%%%%%%%%%%%%%%%%%%%%%%%
\section{INTRODUCTION}

Unmanned aerial vehicles (UAVs) are increasingly used in aerial photography, surveying and mapping, and logistics delivery, where autonomous navigation is essential for safe and efficient operation \cite{xiao2025vision}. 
Conventional UAV navigation often follows a sequential pipeline of perception, mapping, planning, and control, whose strong module coupling can cause error accumulation and substantial onboard computational burden \cite{wang2023learning,zhou2020ego}. 
Learning-based methods offer a more reactive alternative by directly mapping onboard observations to control commands \cite{loquercio2021learning,nguyen2022motion}. 
In particular, reinforcement learning trained in high-fidelity simulation has shown strong performance in complex environments \cite{song2023reaching,kaufmann2023champion}.

Most existing studies on UAV obstacle avoidance characterize scene complexity mainly by the number of obstacles and their spatial density, while paying much less attention to obstacle scale \cite{bhattacharya2025vision,zhao2024learning}. 
However, obstacle scale is not a minor environmental detail, because it directly affects both perception quality and avoidance behavior \cite{chen2025general}. 
Large obstacles can severely occlude free space under the UAV's limited field of view, whereas small obstacles may occupy only a few pixels and are therefore easy to miss \cite{lee2025quadrotor,xu2025flying}. 
As a result, methods that perform well in nominal-scale scenes may still struggle in multi-scale environments where large and small obstacles coexist.

These observations suggest that robust navigation in multi-scale environments requires a policy that can simultaneously preserve small-obstacle sensitivity and handle occlusions caused by large obstacles \cite{kulkarni2023semantically,yu2024mavrl}.
On the one hand, the visual representation should preserve collision-relevant details so that small obstacles remain salient for decision making \cite{kulkarni2023semantically}.
On the other hand, the policy should exploit temporal context to infer scene structure beyond the current field of view and maintain reliable detouring decisions around large occluders \cite{yu2024mavrl}.
Therefore, an effective solution should jointly improve small-obstacle sensitivity and occlusion handling.

Guided by this insight, we propose MeCaRL, a memory-enhanced and collision-aware reinforcement learning approach for UAV obstacle avoidance. 
MeCaRL combines collision-aware representation learning to improve safety around small obstacles with temporal memory to alleviate occlusions caused by large obstacles. 
As illustrated in Fig.~\ref{fig:1}(a), MeCaRL produces safer and more reliable trajectories than existing learning-based baselines in both extra-large and ultra-small obstacle settings. 
Fig.~\ref{fig:1}(b) further shows that MeCaRL achieves higher success rates in these two representative scenarios. 

Our main contributions are as follows:
\begin{itemize}
  \item We introduce multi-scale obstacle environments for UAV visual navigation, including ultra-small and extra-large obstacle regimes, to explicitly evaluate the impact of obstacle scale.
  \item We develop MeCaRL, a reinforcement learning framework that combines collision-aware representation learning for small-obstacle safety with temporal memory for large-obstacle occlusion handling.
  \item We conduct systematic simulation and real-world experiments to evaluate MeCaRL across obstacle scales, demonstrating its effectiveness in both conventional and challenging multi-scale environments.
\end{itemize}


\section{RELATED WORK}

\subsection{Learning-based Visual Navigation}
Learning-based visual navigation is commonly studied through imitation learning and reinforcement learning. 
Imitation learning directly maps visual observations to control commands using expert demonstrations, but its performance is often limited by the expert policy and may generalize poorly in complex environments \cite{loquercio2021learning,bhattacharya2025vision}. 
Reinforcement learning instead optimizes long-term returns through interaction and can better balance safety and efficiency, although end-to-end visual RL often suffers from sample inefficiency and instability \cite{kaufmann2023champion}.

To improve training efficiency and robustness, many recent methods adopt a two-stage pipeline that first learns compact visual representations and then trains policies in the latent space. 
Hoeller et al.~\cite{hoeller2021learning} fuse sequential depth images with camera trajectories to learn local structural representations for dynamic obstacle avoidance.
Kulkarni et al.~\cite{kulkarni2024reinforcement} construct low-dimensional representations from depth-based collision cues, enabling efficient local navigation while accounting for body size.
Yu et al.~\cite{yu2024mavrl} incorporate depth history into the latent space to improve speed adaptation and memory, and further align simulated and real depth representations using a variational autoencoder for cross-domain robustness \cite{yu2025depth}.
For RGB navigation, Zhang et al.~\cite{zhang2025learning} learn depth-consistent RGB representations through cross-modal contrastive learning. 
Despite these advances, most existing methods assume obstacles of similar scales and rarely study how obstacle scale variations affect perception and policy learning.

\subsection{Obstacle Avoidance in Multi-Scale Environments}
In multi-scale environments, existing methods often address small and large obstacles separately. 
Small or slender obstacles are easily missed in visual or depth observations, motivating methods that enhance geometric representations or safety modeling for more reliable avoidance. 
For example, Kulkarni et al.~\cite{kulkarni2023semantically} learn a semantically enhanced depth representation that preserves slender structures under compression,
while Ren et al.~\cite{ren2025safety} and Xu et al.~\cite{xu2025flying} improve avoidance of tiny obstacles and narrow passages using fine-grained LiDAR-based planning or reinforcement learning.

Large obstacles mainly introduce severe occlusions and detouring difficulty, which motivates approaches based on path generation or privileged information. 
Chen et al.~\cite{chen2025general} improve path reliability near large obstacles through local sub-goals and a backtracking detour mechanism.
Lee et al.~\cite{lee2025quadrotor} use privileged cues to maintain the correct detouring direction under occlusion.
Zhai et al.~\cite{zhai2025pa} propose perception-aware MPPI to encourage exploration when the goal is occluded. 
However, jointly handling small-obstacle sensitivity and large-obstacle occlusion under coexisting multi-scale obstacles remains underexplored.

\section{MEMORY-ENHANCED COLLISION-AWARE REPRESENTATION LEARNING}

\begin{figure*}[t]
  \centering
  \includegraphics[width=\textwidth]{figures/fig2_architecture.pdf}
  \caption{Overview of the MeCaRL architecture and training pipeline. 
  A depth image is first encoded into a latent vector, which is then fused with temporal history by an LSTM to produce a memory-enhanced representation for policy prediction. 
  The dashed box denotes the training-only branch, where collision-aware reconstruction is used to learn these representations.}
  \label{fig:architecture}
\end{figure*}

As illustrated in Fig.~\ref{fig:architecture}, the framework encodes each depth image with a VAE and fuses temporal observations with an LSTM to obtain a memory-enhanced representation, 
which is concatenated with the UAV state and target information for policy prediction. 
The training pipeline consists of three stages: collision-aware representation learning, memory-enhanced representation learning, and reinforcement learning for navigation. 
We first train an initial policy to collect depth-image sequences, then train the VAE and LSTM to learn collision-aware and memory-enhanced representations, and finally freeze the representation modules and retrain the policy network for navigation.


\subsection{Dataset Collection and Collision-aware Processing}

\begin{figure*}[t]
  \centering
  \includegraphics[scale=0.2]{figures/fig_ca_preprocess.pdf}
  \caption{Collision-aware preprocessing pipeline. Obstacle boundaries are inflated according to the UAV body size to obtain an edge-inflated depth map, 
  while obstacle interiors are approximated by distance correction. The two results are fused to produce the final collision-aware depth image.}
  \label{fig:3_ca_processing}
\end{figure*}

To build the training dataset, we first train an initial policy while keeping the VAE and LSTM randomly initialized and fixed. 
This policy mainly learns basic goal-reaching behavior and is used to collect depth-image sequences in simulated environments with obstacles of different scales. 
Because the dataset is gathered from continuous trajectories, the resulting observations naturally contain temporal correlations, which are useful for subsequent LSTM training.

After obtaining the raw depth-image dataset, we perform collision-aware preprocessing to account for the UAV body size and generate collision-aware depth images, as illustrated in Fig.~\ref{fig:3_ca_processing}. 
Given a raw depth image $I^{\mathrm{raw}}$, we first extract obstacle contours and back-project them into 3D. 
The contour regions are then inflated according to the effective UAV size $d_{\mathrm{UAV}}$, and rendered to produce an edge-inflation image $I^{\mathrm{infl}}$, which emphasizes collision risks near obstacle boundaries. 
For obstacle interiors, we apply distance correction by subtracting half of the UAV size from the raw depth image, yielding $I^{\mathrm{corr}} = I^{\mathrm{raw}} - d_{\mathrm{UAV}}/2$. 
The final collision-aware depth image is then obtained by pixel-wise fusion, $I^{\mathrm{ca}} = \min\left(I^{\mathrm{infl}}, I^{\mathrm{corr}}\right)$. 
This result provides collision-relevant supervision for the subsequent representation learning stage.

\subsection{Collision-aware Representation Learning}
After obtaining paired raw depth images and collision-aware depth images, we train a VAE to learn collision-aware latent representations from raw depth observations. 
During training, the encoder maps a raw depth image $I_t^{\mathrm{raw}}$ to the parameters of a latent distribution, from which a latent vector $z_t$ is sampled via the reparameterization trick:
\begin{equation}
z_t \sim \mathcal{N}\!\left( \mu_t,\; \mathrm{diag}(\sigma_t^2) \right).
\end{equation}

The decoder reconstructs the corresponding collision-aware image $\hat{I}_t^{\mathrm{ca}}$ from $z_t$. 
In this way, collision-aware images serve as supervision targets for representation learning, encouraging the latent space to preserve geometric structure and safety-related cues that are important for collision reasoning. 
This design embeds collision-aware information into the latent representation without requiring explicit collision-aware preprocessing at inference time.

The training objective of the model consists of a collision-aware reconstruction loss and a latent space regularization term, and is defined as:
\begin{equation}
\begin{aligned}
\mathcal{L}_{\mathrm{CA}} &= \mathcal{L}_{\mathrm{coll}} + \lambda_{\mathrm{KL}}\,\mathcal{L}_{\mathrm{KL}},\\
\mathcal{L}_{\mathrm{coll}} &= \mathrm{MSE}\!\left(I_t^{\mathrm{ca}},\,\hat{I}_t^{\mathrm{ca}}\right),\\
\mathcal{L}_{\mathrm{KL}} &= \frac{1}{2}\sum_{i=1}^{n_z}
\left(1-\mu_{t,i}^2-\sigma_{t,i}^2+\log(\sigma_{t,i}^2)\right),
\end{aligned}
\label{eq:ca_loss}
\end{equation}

where $\lambda_{\mathrm{KL}}$ is the weight of the KL regularization term, $n_z$ denotes the latent dimensionality, 
and $I_t^{\mathrm{ca}}$ and $\hat{I}_t^{\mathrm{ca}}$ denote the ground-truth and reconstructed collision-aware images, respectively.


\subsection{Memory-enhanced Representation Learning}
Using the pretrained encoder, each raw depth observation $I_t^{\mathrm{raw}}$ is first mapped to a latent vector $z_t$. 
The sequence of latent vectors is then processed by an LSTM to obtain a hidden state $h_t$ that summarizes temporal context. 
A prediction head maps $h_t$ to two latent codes, $\tilde z_{t-T}$ and $\tilde z_t$, which are decoded by the frozen VAE decoder into $\hat{I}_{t-T}^{\mathrm{ca}}$ and $\hat{I}_{t}^{\mathrm{ca}}$, respectively. 
By jointly reconstructing past and current collision-aware observations, the supervision encourages $h_t$ to retain temporally consistent scene structure and collision-relevant information, 
which is especially useful when the current observation is partially occluded by large obstacles.

The LSTM is trained by minimizing the sum of reconstruction errors over the two time steps:
\begin{equation}
\mathcal{L}_{\mathrm{LSTM}}=
\mathrm{MSE}\!\left(I_{t}^{\mathrm{ca}},\,\hat{I}_{t}^{\mathrm{ca}}\right)
+
\mathrm{MSE}\!\left(I_{t-T}^{\mathrm{ca}},\,\hat{I}_{t-T}^{\mathrm{ca}}\right).
\end{equation}
Here, $T$ denotes the temporal interval. Minimizing this objective yields a memory-enhanced representation $h_t$ for downstream policy learning.

\section{REINFORCEMENT LEARNING FOR NAVIGATION}
\subsection{Task Formulation}

We formulate UAV navigation as a partially observable Markov decision process (POMDP). The policy
$\pi_\theta(a_t\mid o_t)$ is trained to maximize the expected discounted return under partial observations.
The action is defined as a high-level control command $a_t=[a_x,a_y,a_z,v_{\mathrm{yaw}}]$, where
$[a_x,a_y,a_z]$ denotes the desired acceleration and $v_{\mathrm{yaw}}$ denotes the yaw rate.
The observation is formed by concatenating the UAV state and target information $s_t$ with the
memory-enhanced representation $h_t$, i.e., $o_t=[s_t,h_t]$. Specifically,
\begin{equation}
s_t=\big[\log d_{\mathrm{hor}},\ d_z,\ d_{\mathrm{norm}},\ v^{\mathrm{W}},\ e,\ \omega^{\mathrm{B}}\big].
\end{equation}
Here, $d_{\mathrm{hor}}$ and $d_z$ denote the horizontal and vertical distances to the target,
$d_{\mathrm{norm}}$ is the normalized target direction, $v^{\mathrm{W}}$ is the linear velocity in the
world frame, $e$ denotes the Euler-angle attitude, and $\omega^{\mathrm{B}}$ is the body-frame angular velocity.

% The perceptual representation $h_t$ is obtained by encoding
% the depth observation $I_t$ and fusing historical information with the LSTM, providing a memory-enhanced
% collision-aware representation. With the above observation and action definitions, the policy network takes $o_t$
% as input and outputs $a_t$, and the policy parameters are optimized via reinforcement learning to obtain stable
% collision-free navigation behaviors.

\subsection{Reward Function}
We design the reward as the combination of terminal rewards and a dense progress reward. An episode terminates when the UAV reaches the goal, collides with obstacles, or leaves the allowed flight region. The overall reward at time step $t$ is defined as
\begin{equation}
r_t=
\begin{cases}
r_{\mathrm{arrive}}, & \|d_t\|<d_{\min},\\
r_{\mathrm{collision}}, & \text{a collision occurs},\\
r_{\mathrm{exceed}}, & p_t \notin [p_{\min},\,p_{\max}],\\
r_{\mathrm{prog}}, & \text{otherwise}.
\end{cases}
\end{equation}
Here, $d_t$ denotes the relative displacement to the goal and $p_t=[x_t,y_t,z_t]$ is the UAV position. $d_{\min}$ is the arrival threshold, and $[p_{\min},p_{\max}]$ denotes the valid flight region defined by the lower and upper bounds along each axis.

The progress reward is defined as
\begin{equation}
\begin{aligned}
r_{\mathrm{prog}}=&\ \lambda_d \log d_{\mathrm{hor}} + \lambda_z d_z
+ \lambda_v\,\max(0, v_{\mathrm{hor}}-v_{\max})\,v_{\mathrm{hor}} \\
&+ \lambda_{\mathrm{dir}}\left\|d_{\mathrm{norm}}-v_{\mathrm{norm}}\right\|_1
+ \lambda_{\mathrm{ang}}\left\|\omega^{\mathrm{B}}\right\| \\
&+ \lambda_{\mathrm{lat}}\left(|v_y^{\mathrm{B}}|+\max(0,-v_x^{\mathrm{B}})\right),
\end{aligned}
\end{equation}
where $d_{\mathrm{hor}}$ and $d_z$ are the horizontal and vertical distances to the goal,
$v_{\mathrm{hor}}$ is the horizontal speed with limit $v_{\max}$, $d_{\mathrm{norm}}$ and
$v_{\mathrm{norm}}$ are the normalized goal and velocity directions, $\omega^{\mathrm{B}}$ is the body-frame
angular velocity, and $v_x^{\mathrm{B}}$, $v_y^{\mathrm{B}}$ are the forward and lateral body-frame velocities.
The distance and direction terms encourage goal reaching, while the remaining terms regularize
speed, angular motion, and lateral/backward movements for stable and safe flight.

\section{EXPERIMENTS}

\subsection{Simulation Setup} 

\begin{figure*}[t]
  \centering
  \includegraphics[width=0.95\textwidth]{figures/fig_sim_env.pdf}
  \caption{Simulation environments with different obstacle scales. (a) Nominal-scale training environment. 
  (b)--(g) Test environments with ultra-small (US), small (S), medium (M), large (L), extra-large (XL), and mixed-scale (MIX) obstacles.}
  \label{fig:sim_env}
\end{figure*}

We build the simulation environment in Unity, following the reinforcement-learning setup of MAVRL \cite{yu2024mavrl} with a simplified dynamics model. 
To evaluate robustness to obstacle-scale variation, the policy is trained in a nominal-scale environment and then tested in six environments with different obstacle sizes, as shown in Fig.~\ref{fig:sim_env}.

The scene is an indoor garage populated with cylindrical obstacles. We construct seven environment configurations by varying obstacle diameter: 
one nominal-scale training environment (10--50 cm) and six test environments, namely ultra-small (US, 1--5 cm), small (S, 10--30 cm), medium (M, 40--80 cm), large (L, 100--200 cm), extra-large (XL, 400--500 cm), 
and mixed-scale (MIX, 1--500 cm). Obstacle density is controlled by the radius parameter used for Poisson obstacle placement, yielding different clutter levels and navigation difficulty.

For representation learning, we collect trajectory data in the nominal-scale training environment using the initial exploration policy. 
The resulting dataset contains 200 episodes and approximately 10,000 depth images for subsequent representation learning and model training.

% \subsection{Analysis of Latent Representations}
\subsection{Representation Ablation Study}

\begin{table}[t]
  \centering
  \caption{EFFECT OF TEMPORAL INTERVAL $T$}
  \label{tab:t_sensitivity}
  \small
  \setlength{\tabcolsep}{4pt}
  \renewcommand{\arraystretch}{1.05}
  \begin{tabular}{c cccccc}
    \toprule
    $T$ & 1 & 5 & 10 & 15 & 20 & 30 \\
    \midrule
    Success rate & 0.64 & 0.93 & \textbf{0.96} & 0.87 & 0.82 & 0.74 \\
    \bottomrule
  \end{tabular}
\end{table}

We first analyze the effect of the temporal interval $T$ in the memory module under mixed-scale obstacle settings. 
Table~\ref{tab:t_sensitivity} shows that the success rate increases from 0.64 at $T=1$ to 0.96 at $T=10$, and then decreases as $T$ becomes larger. 
This trend indicates that a moderate temporal interval provides useful historical context, whereas excessively short or long intervals are less effective. 
Based on this sensitivity analysis, we use $T=10$ in the following experiments.


\begin{table*}[t]
  \centering
  \caption{REPRESENTATION ABLATION ACROSS OBSTACLE SCALES}
  \label{tab:repr_eval}
  \small
  \setlength{\tabcolsep}{5pt}
  \renewcommand{\arraystretch}{1.15}
  \begin{tabular}{c cccc cccc}
    \toprule
    \multirow{2}{*}{Environment} &
    \multicolumn{4}{c}{Success rate} &
    \multicolumn{4}{c}{Average speed (m/s)} \\
    \cmidrule(lr){2-5}\cmidrule(lr){6-9}
    & VanillaRL & MeRL & CaRL & MeCaRL
    & VanillaRL & MeRL & CaRL & MeCaRL \\
    \midrule
    Ultra-small & 0.25 & 0.42 & 0.65 & \textbf{0.91} & 1.74 & 1.25 & \textbf{2.35} & 2.33 \\
    Small       & 0.89 & \textbf{1.00} & \textbf{1.00} & \textbf{1.00} & 2.07 & 1.40 & 2.06 & \textbf{2.25} \\
    Medium      & 0.31 & 0.96 & 0.98 & \textbf{1.00} & 1.69 & 1.29 & 2.00 & \textbf{2.33} \\
    Large       & 0.48 & 0.91 & 0.92 & \textbf{0.96} & 1.58 & 1.26 & \textbf{2.46} & 2.29 \\
    Extra-large & 0.39 & 0.62 & 0.60 & \textbf{0.74} & 1.35 & 1.28 & \textbf{1.93} & 1.72 \\
    Mixed       & 0.08 & 0.54 & 0.91 & \textbf{0.96} & 0.88 & 1.17 & 1.70 & \textbf{1.86} \\
    \bottomrule
  \end{tabular}
\end{table*}

To evaluate the contribution of memory enhancement and collision awareness, we compare four policy variants under the same reinforcement-learning framework: VanillaRL, MeRL, CaRL, and MeCaRL. 
VanillaRL uses neither module, MeRL uses memory enhancement only, CaRL uses collision awareness only, and MeCaRL combines both. 
All variants are evaluated in densely cluttered environments with randomly placed obstacles. Each configuration is tested with four random seeds and 25 runs per seed, resulting in 100 runs in total.

As shown in Table~\ref{tab:repr_eval}, VanillaRL performs poorly in densely cluttered multi-scale environments, especially in the ultra-small, extra-large, and mixed settings. 
Adding memory enhancement improves success rates across scales, with more evident gains in the medium and large environments, suggesting that temporal context helps mitigate partial observability caused by occlusions. 
Adding collision awareness yields larger gains in the ultra-small and mixed environments, indicating improved sensitivity to fine structures and safety margins. 
MeCaRL achieves the highest success rate in all environments while maintaining high flight speed in most cases. 
These results suggest that collision awareness and memory enhancement play complementary roles in multi-scale obstacle avoidance.

\subsection{Baseline Comparisons}

% figure*和table*只支持t和b，不支持h
\begin{figure*}[t]
  \centering
  \includegraphics[width=0.70\textwidth]{figures/fig_trajectory.pdf}
  % \includegraphics[width=0.80\textwidth]{figures/fig_trajectory.pdf}
  \caption{Representative multi-run trajectories of Agile-autonomy, MAVRL, and MeCaRL in extreme-scale environments. 
  Top row: ultra-small. Bottom row: extra-large. Red circles indicate failed trajectories, and trajectory colors denote velocity.}
  \label{fig:traj}
\end{figure*}

\begin{table*}[t]
  \centering
  \caption{BASELINE COMPARISON ACROSS OBSTACLE SCALES}
  \label{tab:baseline_eval}
  \small
  \setlength{\tabcolsep}{5pt}
  \renewcommand{\arraystretch}{1.10}
  \begin{tabular}{c ccc ccc}
    \toprule
    \multirow{2}{*}{Environment} &
    \multicolumn{3}{c}{Success rate} &
    \multicolumn{3}{c}{Average speed (m/s)} \\
    \cmidrule(lr){2-4}\cmidrule(lr){5-7}
    & Agile-autonomy \cite{loquercio2021learning} & MAVRL \cite{yu2024mavrl} & MeCaRL
    & Agile-autonomy \cite{loquercio2021learning} & MAVRL \cite{yu2024mavrl} & MeCaRL \\
    \midrule
    Ultra-small & 0.27 & 0.29 & \textbf{0.77} & 1.67 & 1.22 & \textbf{2.08} \\
    Small       & 0.32 & 0.46 & \textbf{0.83} & 1.72 & 1.34 & \textbf{2.04} \\
    Medium      & 0.34 & 0.71 & \textbf{0.91} & 1.70 & 1.32 & \textbf{1.87} \\
    Large       & 0.26 & 0.46 & \textbf{0.90} & 1.68 & 1.20 & \textbf{2.05} \\
    Extra-large & 0.24 & 0.44 & \textbf{0.72} & 1.64 & 1.18 & \textbf{1.68} \\
    Mixed       & 0.32 & 0.56 & \textbf{0.79} & 1.69 & 1.21 & \textbf{1.89} \\
    \bottomrule
  \end{tabular}
\end{table*}


To evaluate MeCaRL against existing methods, we conduct benchmark comparisons in AvoidBench \cite{yu2023avoidbench}, which is built on Flightmare \cite{song2021flightmare} for vision-based obstacle avoidance. 
Following the same environment setup as described above, we evaluate all methods in six test environments with varying obstacle diameters under densely cluttered and randomly generated layouts. 
The compared baselines include MAVRL \cite{yu2024mavrl}, a reinforcement-learning-based method, and Agile-autonomy \cite{loquercio2021learning}, a supervised-learning-based method. 
For MeCaRL, the predicted acceleration commands are converted into executable body-rate and thrust commands using the Agilicious MPC controller \cite{foehn2022agilicious}.
All methods are evaluated under the same observation setting, termination conditions, and environment generation protocol.

Table~\ref{tab:baseline_eval} summarizes the performance across obstacle scales. 
All values in Table~\ref{tab:baseline_eval} are averaged over 100 runs. 
MeCaRL achieves the highest success rate in all six environments and also maintains competitive average speed, indicating a favorable trade-off between reliability and efficiency. 
The gains are particularly clear in the ultra-small and extra-large settings, which represent two extreme-scale navigation conditions. 
In the ultra-small environment, MeCaRL improves the success rate from 0.29 to 0.77 relative to MAVRL, suggesting better sensitivity to tiny obstacles. 
In the extra-large environment, MeCaRL improves the success rate from 0.44 to 0.72 while maintaining higher average speed, indicating stronger robustness under severe occlusions.

Fig.~\ref{fig:traj} shows representative multi-run trajectories in the ultra-small and extra-large environments. 
In the ultra-small setting, Agile-autonomy and MAVRL exhibit more failed trajectories and less consistent rollouts toward the goal, whereas MeCaRL produces more coherent trajectories with fewer failures. 
In the extra-large setting, MeCaRL also shows more reliable detouring behavior around large obstacles, while Agile-autonomy and MAVRL exhibit more failures and less efficient path selection. 
These trajectory patterns are consistent with the quantitative results in Table~\ref{tab:baseline_eval}.

\subsection{Network Attention Visualization}

To qualitatively examine the visual regions emphasized by the learned policy, we use Grad-CAM \cite{selvaraju2017grad} to visualize the contribution of different spatial locations in the input depth image. 
Warm colors indicate stronger influence on the policy output, while cool colors indicate weaker influence.

\begin{figure*}[t]
  \centering
  \includegraphics[width=0.85\textwidth]{figures/fig_attention_visualization.pdf}
  \caption{Grad-CAM attention visualizations of MAVRL and MeCaRL across obstacle scales. Warm colors indicate stronger influence on the policy output.}
  \label{fig:attention_vis}
\end{figure*}

\begin{figure*}[t]
  \centering
  \includegraphics[width=0.95\textwidth]{figures/fig_outdoor_exp.pdf}
  \caption{Real-world outdoor experiments. (a) The deployed quadrotor platform. (b) A representative flight trajectory from the start to the goal, where the color indicates the flight speed. (c) The outdoor test environment.}
  \label{fig:outdoor_exp}
\end{figure*}

We apply Grad-CAM to scenes with obstacles at different scales and compare MeCaRL with MAVRL, as shown in Fig.~\ref{fig:attention_vis}. 
In scenes containing small-scale and medium-scale obstacles, MeCaRL tends to assign stronger attention to small obstacles and potential risk regions at longer distances, 
whereas MAVRL often exhibits weaker or more diffuse responses to these structures. 
In scenes with large-scale obstacles, MeCaRL shows more concentrated responses on obstacle bodies and contours, while MAVRL distributes attention more broadly over both obstacle and background regions. 
MeCaRL also tends to highlight surrounding boundary regions near obstacles, suggesting improved sensitivity to safety margins. 
These qualitative patterns are consistent with the superior navigation performance of MeCaRL in multi-scale environments.


\subsection{Real-World Tests}
% \begin{figure*}[!t]
%   \centering
%   \includegraphics[width=0.95\textwidth]{figures/fig_outdoor_exp.pdf}
%   \caption{Real-world outdoor experiments. (a) The deployed quadrotor platform. (b) A representative flight trajectory from the start to the goal, where the color indicates the flight speed. (c) The outdoor test environment.}
%   \label{fig:outdoor_exp}
% \end{figure*}

To examine the deployability of MeCaRL in outdoor scenarios, we deploy the method on a quadrotor with a 250 mm wheelbase, as shown in Fig.~\ref{fig:outdoor_exp}(a). 
The platform is equipped with an Intel RealSense D435 depth camera, an NVIDIA Jetson Orin NX onboard computer, and a Pixhawk 6C mini flight controller for onboard sensing, computation, and control.

To reduce the sim-to-real gap, we fine-tune the VAE and LSTM representation modules using depth data collected in the real environment, while keeping the policy network trained in simulation. 
The network architecture and hyperparameter settings are otherwise kept consistent with those used in simulation.

We conduct outdoor flight tests in a tree-rich environment, shown in Fig.~\ref{fig:outdoor_exp}(c), where the quadrotor is required to reach a target while avoiding obstacles. 
Fig.~\ref{fig:outdoor_exp}(b) shows a representative real-world flight trajectory. 
These results provide qualitative evidence that MeCaRL can be deployed for outdoor obstacle avoidance and goal-reaching with the onboard sensing and computation setup.


\section{CONCLUSION}
This paper presents MeCaRL, a memory-enhanced and collision-aware reinforcement learning framework for UAV navigation in multi-scale obstacle environments. 
By combining collision-aware representation learning from depth observations with temporal memory modeling, MeCaRL improves obstacle avoidance performance across obstacle scales, 
especially in challenging ultra-small and extra-large settings. Simulation results show that MeCaRL consistently outperforms baseline methods in success rate while maintaining competitive flight efficiency in multi-scale environments. 
Real-world outdoor flight tests further provide qualitative evidence that MeCaRL can be deployed for obstacle avoidance and goal-reaching with onboard sensing and computation. 
Future work will extend the framework to dynamic obstacles and more complex real-world environments.



% Balance the two columns on the final page.
% 用来调整最后一页的内容大小，防止参考文献在双栏显示左栏比右栏长太多，不美观
\balance


\bibliographystyle{IEEEtran}
\bibliography{mecarl}

\end{document}
